\chapter{Introducción}

\section{Descripción del problema}

Las actas del Pleno del Ayuntamiento de Bilbao son el registro oficial de todas las
proposiciones, debates y votaciones municipales desde 2007 hasta la actualidad. Se
publican como documentos PDF, uno por sesión, y cada acta puede superar el centenar
de páginas. El corpus completo utilizado en este trabajo reúne cerca de veinte años de
actas y más de 3.900 proposiciones distintas.

Este formato tiene un problema práctico: encontrar información concreta --qué votó
un grupo político sobre un tema, cuántas proposiciones sobre vivienda se presentaron en
un año determinado, cómo ha evolucionado un asunto a lo largo de los mandatos-- exige
leer manualmente decenas de documentos extensos, o como mínimo saber de antemano en
qué fecha y página buscar. No existe ninguna vía para hacer preguntas en lenguaje natural
sobre este archivo y obtener una respuesta verificable, con cifras y citas a la fuente
original.

\subsection{Acceso a la información municipal}

El objetivo de este TFG es cerrar esa brecha: construir un sistema que permita consultar
en castellano, mediante preguntas formuladas de forma natural, el contenido íntegro de
las actas del Pleno, y que responda con datos verificables --enlazados a la página exacta
del PDF de origen-- en lugar de con texto generado sin respaldo documental.

\section{Estructura de la memoria}

Esta memoria se organiza como sigue. El Capítulo~\ref{ch:contexto} introduce los
conceptos técnicos en los que se apoya el sistema (modelos de lenguaje, RAG,
grafos de conocimiento y GraphRAG). El Capítulo~\ref{ch:objetivos} concreta los
objetivos del trabajo y su alcance real. El Capítulo~\ref{ch:beneficios} discute
los beneficios del sistema y su relación con los Objetivos de Desarrollo Sostenible.
Los capítulos siguientes describen el estado del arte, el análisis de alternativas
tecnológicas, los riesgos del proyecto y, con más detalle, la solución desarrollada:
los dos sistemas de recuperación de información construidos --uno basado en
búsqueda vectorial y otro en un grafo de conocimiento RDF/OWL consultado por
SPARQL-- y la interfaz web conversacional que los expone al usuario. La memoria
cierra con la planificación seguida, el presupuesto, las conclusiones y los anexos.
